
\subsection{Dependent Measure}

On each trial, participants produce an adjustment that may differ from the ground truth. We scored the illusion effect as follows: For the Brentano Illusion, the score was the number of pixels to the left (or right) of half of the horizontal segment relative to the size (in pixels) of the segment. Relative effects were used for the Ebbinghaus and Ponzo illusions, each effect was the difference between the adjusted value and the true value relative to the size of the true value. For the Poggendorff and Zöllner illusions, however, there is nothing to compare the adjustments against. Therefore, the scores were measured as the number of displaced pixels (Poggendorff) from the true value or the angle from horizontal (Zöllner).  We call these values the "adjustment score" in the following and they are the relative bias in adjustment.



\subsection{Data Cleaning}

Data obtained online can be messy.  Surprisingly, the data in this report were fairly clean, and we think this is because the battery was brief (it took under 20 minutes to complete) and participants reported that it was engaging.  Because we had little expectations on the ranges of data {\it a priori}, we performed data cleaning after data collection but before analysis of pattern of correlations.   A very detailed report on data cleaning is available at \href{https://github.com/specl/ind-illusions/blob/public/_data-processing/data-cleaning.pdf}{linked document}\footnote{\url{https://github.com/specl/ind-illusions/blob/public/_data-processing/data-cleaning.pdf}}; here is a summary: Data-cleaning decisions fall into three camps: those made without any reference to the adjustment score, those made with reference to adjustment scores in each task in isolation, and those made with reference to adjustment scores across two tasks at a time.  


{\bf Data Cleaning Without Reference To Adjustment Sores.} We first made a set of cleaning decisions before any analysis of the adjustment score:  I. One participant performed the task twice rather than once.  The second session was excluded.  II.  We worried about participants who sped through the task at atypical and unreasonably fast speed.  As a bit of context, the adjustment paradigm here consisted of making several keypresses to finely adjust objects. This process takes longer than conventional response times, and typical responses are in the 5-second range.  The histograms of response times was bimodal with a separate mode below 600 ms. This mode was not large, and overall, only 1.5\% of responses fell here. These responses were excluded, as were individuals who had more than 5\% of their responses in this mode (5 individuals). 

{\bf Data Cleaning Each Task In Isolation.} We worried about participants who may not have understood some of the adjustment tasks.   We conducted a univariate analysis to detect outliers based on standardized scores across different tasks and their respective versions. Participants whose z-score on any task exceeded 4 were removed (6 individuals).   When these z-scores exceeded 4, they often did so dramatically indicating that these 6 individuals did not understand the instructions for one or more of the tasks.


{\bf Data Cleaning Across Bivariate Configurations.} We worried that correlations across tasks might be unduly influenced by high-leverage points.  We ran a suite of leverage tests (Difference in Fits, Difference in Beta Coefficients, and Cook’s Distance \parencite{Belsley.etal.1980}).  Although there were a few high leverage points, they were not unreasonable in any other way and we could find no rationale for excluding them.

\subsection{Demographics}

After the data-cleaning process, 138 participants remained. The average age was 35.34 years (SD = 11.01). There were 66 females, 71 males, and 1 participant who preferred not to say.

\subsection{Tasks In Isolation}


The goal of the analysis is to understand the relationships among the tasks. However, tasks must be suitable in that they measure a signal without an undue degree of trial noise. If individual task scores are too noisy, then correlations are attenuated too low, and latent-variable decompositions may be unstable \parencite[]{Karr.etal.2018,Rouder.etal.2023}. The first step, then, is to provide a deep dive into the trial noise in the task scores themselves.


{\bf Size of the effects.}  When individual-difference studies are conducted, the focus is on covariation and correlation, and the size of the effects does not enter into the critical parts of the analyses.  Yet, the effect sizes may serve as an indicator of the robustness of subsequent individual-difference analyses.  When effects are large, there is substantial headroom for people to vary. When effects are small, there is far less headroom. A good first check, then, is the size of the effect

We computed an effect size and a $t$-value for every individual in each of the 10 tasks.  The effect size for the $i$th individual in the $j$th task is $d_{ij}=m_{ij}/s_{ij}$, where $m_{ij}$ and $s_{ij}$ are respectively the sample mean and sample standard deviation of adjustment scores for all trials for the $i$th individual in the $j$th task.  The $t$-value is $t_{ij}=d_{ij}\times\sqrt{K_{ij}}$.  Figure~\ref{fig:effectPlots}A shows a histogram of all effect sizes and $t$-values, and they are nothing short of spectacular in value.  Even though there are at most 15 trials per individual per task version, the vast majority of the effect sizes exceed Cohen's 0.8 criteria for large, and even more are significant at 0.05 level.   Figure~\ref{fig:effectPlots}B shows these effects tabulated by task as empirical cumulative distribution functions.  Even though there are differences across the tasks in mean effect sizes and the spread of effect sizes, the large effect sizes hold for all 10 tasks.  In summary, visual illusions are so large that performance from just 15 trials contains more signal than typical cognitive studies comprised of thousands of trials.

{\bf Discriminating among individuals within a task.}  For the remainder of the analyses, we used scaled adjustment scores as follows.  For each task, mean squared error (MSE) was computed among individuals, and the square-root is the standard deviation of task-specific residual variance.  Then, each score for all individuals in a task was expressed in units of this standard deviation, that is, in task-specific effect size units.  Let $Y_{ijk}$ be the $k$th replicate observation for the $i$th person on the $j$th task.  Let $s_j=\sqrt{\mbox{MSE}_j}=\left(\sum_{ik}(Y_{ijk}-m_{ij})^2/I(K-1)\right)^{1/2}$. Then the task-scaled score is $Y_{ijk}/s_j$.  Figure~\ref{fig:taskPlots} shows these scaled scores in a configuration that we term {\em task plots}. Here, individuals' observed mean scaled scores in each task are ordered and plotted using dashed lines.  The solid lines and shaded areas reflect regularized estimates (from a Bayesian hierarchical model discussed subsequently) and 95\% credible intervals, and these intervals show how well each individual's score may be localized.  There are substantial differences in scores across individuals in all tasks, indicating that performance can be discriminated to high degree across individuals. We recommend that task plots always be included as they provide a clear visualization of what is possible in terms of identifying and understanding individual differences.


{\bf Goodness of Task Measures.}  The usual way to understand the goodness of a task is reliability.\footnote{One course is to split the trials, compute the correlation, and scale it up to the full sample size through the Brown-Spearman prophecy formula.  This is a test-theory approach that fails to leverage the fact that experiments are comprised of multiple replicate trials.  Reliability here is computed directly from within-trial and across-trial variabilities, see \textcite{Rouder.Mehrvarz.2024}.}.    There is, of course, a direct relationship.  The larger the spread of the mean scores relative to CIs in the task plots in Figure~\ref{fig:taskPlots}, the higher the reliability.  Yet, both task plots and the associated reliability measure strike us as incomplete for communication about tasks.  



This incompleteness comes from the fact that reliability and task plots are directly influenced by the number of replicate trials.  The Stroop task, for example, is unreliable with 15 trials, but it is very reliable with several thousand trials \parencite[]{Lee.etal.2023}\footnote{c.f. \textcite{Burgoyne.etal.2023} for a high-reliability Stroop-like task where performance is averaged rather than contrasted across congruent and incongruent conditions.}.  What is needed for grounded communication is a task measure of signal-to-noise that does not depend on the number of replicate trials.  We advocate the signal-to-noise ratio of a task as such a measure \parencite[]{Rouder.Mehrvarz.2024}, and denote it as $\gamma^2$.  The signal, the numerator of $\gamma^2$, is the variability across individuals' true scores in a task.  The noise, the denominator of $\gamma^2$, is the trial-to-trial variability. 

The measure $\gamma^2$ may be thought of as a ratio version of ICC(1) \parencite[]{Liljequist.etal.2019}, and we prefer it to ICC because we find the ratio-interpretation quite natural.

 
 Tasks with high values of $\gamma^2$ are reliable even with few replicate trials per individual per task; tasks with low values are reliable only if there are a high number of replicate trials per individual per task.  The $\gamma^2$ for each illusion is denoted in the task plot, and the average is 1.145.  This value means that across illusions, variation across people is as big as variation within a trial.  It is a large value by any standard, and it should be possible to understand individual differences to high precision.

{\bf Split-Half Reliability: A Sanity Check.}  In the preceding plots, we showed that the illusion effects are quite large and that individuals can be discriminated.  We provide one final check on this sanguine state---a split-halves scatter plot.  Plotted are individual mean effects for the second-half of trials as a function of the mean effects for the first-half of trials (Figure~\ref{fig:relPlots}).  As can be seen, the association---the split half reliability---is quite high.  

In summary, with such large and reliable effects, it should be straightforward to assess the relations among tasks without excessive variability or attenuation.

\subsection{Relations Among Tasks}

Now that we have established that these tasks have very high signal-to-noise characteristics, the next step is analyzing the relations among them.  Figure~\ref{fig:corPlot} shows the correlations among the tasks obtained in two different ways.  The upper triangle shows the conventional sample correlations; the lower triangle shows estimates from a hierarchical Wishart model that disattenuates these correlations for trial noise.  The hierarchical Wishart model is from \textcite{Rouder.etal.2023}, and it is comprised of two levels: a data level and an individual-effects level.  At the data level, it is assumed that performance on a trial results is the sum of individual-by-task true score and trial noise.  Let $\theta_{ij}$ denote this true score for the $i$th individual on the $j$th task, and let $\bftheta_i=(\theta_{i1},\ldots,\theta_{iJ})$ denote the vector of true task scores for the $i$th individual.  At the individual-effect level, a multivariate random effects model is placed of $\bftheta_i$: $\bftheta_i \sim \mbox{Normal}_J(\bfmu,\bfSigma)$.  The key here is that all relations among the tasks are in the variance $\bfSigma$.  We estimate $\bfSigma$ in the Bayesian framework, and as such, a prior is needed.  The Wishart prior \parencite[]{Wishart.1928} serves as practical choice as it places little \textit{a priori} constraint on $\bfSigma$ while ensuring that posterior values are valid variance matrix (that is, they are symmetric and positive-definite matrices).  The advantage of the Bayesian framework is two-fold.  First, it is conceptually straightforward approach for separating trial noise from across-individual variation.  Second, analysis is computationally convenient with modern MCMC sampling methods \parencite[]{Gelfand.Smith.1990}.  Posterior of correlations, derived form the variance matrix, are computed within each iteration of the MCMC chains.  The posterior mean is shown in the lower triangle of Figure~\ref{fig:corPlot}.

Comparisons of the sample correlations and model-based correlations reveal a tight correspondence indicating little attenuation in sample correlations.  This lack of attenuation is relatively rare in experimental settings and reflects that effects are large relative to trial noise.  First, the within-version correlations.  For three illusions, the Ponzo, Poggendorf, and Zoellner, there are high correlations across the two versions of the illusions; for the other two illusions, the Brentano and Ebbinghaus, the correlations are more modest.  This result is surprising.   We had expected high correlations, and the task versions were included as a control for response biases.  The modest correlations indicate that it is not wise to average scores across versions.  Second, the cross-illusion correlations.  Each of these is positive in direction and modest in value.  

{\bf Uncertainty In Correlations.}  We recommend that researchers study the uncertainty in cross-task correlations.  In designs where participants provide repeated trials per task score, there are two sources of uncertainty.  One is from the noise across trials.  The degree of uncertainty from this source reflects not only the degree of trial-to-trial variation, but also the number of trials in each cell.  The more trials, the less influence this variation has on the uncertainty of the correlation coefficient.  The other source reflects variation across individuals, and the more individuals the less uncertainty from this source. These sources may be combined into an overall assessment with the hierarchical Wishart model.  The 95\% posterior credible interval (CIs) serves as a suitable measure of uncertainty.  

Figure~\ref{fig:ciCorPlots} shows the mean and uncertainties (95\% CIs) for the correlations.  Consideration of them provides the following insights:  First, as noted, many of them are undoubtedly positive providing evidence for a positive manifold across illusions.  Accordingly, there is at least one factor in common.  Second, the three highest correlations, which stand apart, are correlations among different versions of the tasks.  If we focus on correlations across different illusions, there is not so much variability.  These are clustered around .22 in value with credible intervals that have a full width of .134 ($\pm .067$) on average.  This small value, .22, is by-and-large confirmation of previous results.  One visual illusion accounts for about 4.7\% of the variance in another illusion. 

\subsection{Latent Variable Analyses}

Interpreting sample correlation such as those in Figure \ref{fig:corPlot} is often aided by latent-variable models.   Latent variable modeling, unfortunately, is a bit of a dark art in our view.  While implementation and analysis is relatively straightforward in modern packages, making judicious choices in specification and providing sound interpretation is challenging \parencite[]{Preacher.MacCallum.2003, Armstrong.1967}.  In visual perception, \textcite{Tulver.2019} recently notes that many of these issues, what is a factor, how to measure factors, and how to interpret factors, remain unresolved in the field.  Here, even though the data are of high quality, we run head-long into these difficulties. 



\textbf{Principle Components.}  Many researchers who search for factor structure use principle component analysis (PCA).   Table~\ref{tab:EFA-10} shows the eigenvalues accounted for by each component. Accordingly, there are four factors with eigenvalues greater than 1.0. 

PCA is often seen as interchangeable with factor analysis.  Though similar, there are critical differences.  For experimental psychological research, the latter is superior to the former.  To see the key difference, consider the test case where there is no structure---that is, the correlation matrix is a diagonal matrix plus a small random variation from trial noise.  The PCA for $J$ tasks would yield $J$ factors each accounting for about $1/J$ proportion of the variance.   The factor model underlying factor analysis, however, has error terms that account for uncorrelated noise.  For the test case, the residuals from these error terms would be large, and the factor loadings on illusions that account for the structure would be quite small.  This difference has been well known \parencite[]{Velicer.Jackson.1990} and goes under the moniker that factor analysis partitions common variance rather than total variance \parencite[]{Kline.2023}.  Because factor models are much easier to interpret in low-signal scenarios, they should be used instead of PCA in vision, cognitive control, and related fields.

{\bf Classic Exploratory Factor Analysis.}  We performed exploratory factor analysis (EFA) separately for 1 to 6 factors.  Table~\ref{tab:EFA-10} also shows the model selection statistics, and as can be seen, the situation is ambiguous.  While inspection of eigenvalues (first column) indicates support for a one-factor or four-factor solution, AIC indicates a five-factor solution and BIC indicates a three-factor solution.  Examination of the varimax factor loadings for 5 factors is shown in Table \ref{tab:Fac-10}. Factors tend to load on a single task with only modest cross-task loading.  With this result, as well as the ambiguity of dimensions, it is difficult to gain much additional insight above that from the sample correlations.  

{\bf Classic Confirmatory Factor Analysis.} We also performed a series of confirmatory factor analyses with models shown in Figure~\ref{fig:cfa}.  In the models in Figure~\ref{fig:cfa}A-B, the observed illusion scores  load onto a latent variable.  Correlations across the five latent variables are then computed (Fig~\ref{fig:cfa}A) or modeled with additional latent factors (Fig~\ref{fig:cfa}B).  The main problem we encountered was that the resulting parameters had substantial negative variance components rendering them inadmissible. The problem is a known as a Heyward Condition \parencite[]{Rindskopf.1984} and is under active consideration in the literature \parencite[]{Cooperman.Waller.2022}, We also tried a more flexible version of the model, Figure~\ref{fig:cfa}C), but encountered the same issues.

{\bf Bayesian Hierarchical Factor Models}. We have recently developed hierarchical factor models \parencite[]{Rouder.etal.2024a}. The first level of these models is similar to the hierarchical Wishart model discussed in the previous section. It accounts for trial noise through a linear framework that captures variability within individuals and tasks, allowing for the separation of trial noise from the true-score parameter $\theta_{ij}$.  The second level models these true-score parameters as coming from a latent factor structure:
\[
\pmb{\theta}_i \sim \mathcal{N}_J(\pmb\mu, 
\pmb\Lambda\pmb\Lambda' + D(\pmb\sigma^2)),
\]
where $\pmb\mu$ is the vector of task means, $\pmb\Lambda$ is the factor loading matrix, and $D(\pmb\sigma^2)$ is a diagonal matrix which task-specific unique variances.  The crux of these models is the decomposition of the covariance into common and unique sources of variability across tasks \parencite[]{Preacher.MacCallum.2003}. This decomposition separates shared factor structures from task-specific residual variance. Additionally, factor analysis can take orthogonal or oblique forms depending on whether factor scores are assumed to be independent. The model described here assumes orthogonality, meaning that individuals who score high on one factor are, on average, no better or worse on another. In other words, the model treats abilities across different factors as statistically independent.

We conducted confirmatory factor-model analysis with seven factors \footnote{For full analysis and more detailed model specification, refer to \textcite{Rouder.etal.2024a}.}
. The first five factors were illusion specific, e.g., a Brentano factor, and accounted for correlations across the versions of the task.  The sixth and seventh factors were general in that they could load on several illusions.  These two factors are the main target of study and they are estimated free of influence from across-version variation.  

Table~\ref{tab:CFA-Visual-Illusions} presents the results of the model, and Figure~\ref{fig:EFHBM} visualizes key components of the variance decomposition. The first five factors capture variance specific to differences between versions of each illusion. The key result is that sixth factor loads well on all 10 tasks.  It explains more than four times the variance of the seventh factor. This pattern strongly supports the conclusion that a single general factor underlies susceptibility to visual illusions. A projection of each task-version on the two general factor is depctided by (Figure~\ref{fig:EFHBM}A). However, it is important to note that unique variance between the two versions of specific illusions is relevant (Figure~\ref{fig:EFHBM}B), and unique variability from each version and task  remains substantial (Figure~\ref{fig:EFHBM}C). The general factor alone accounts for approximately 23.3\% of the total variance—substantially higher than the 5\% expected from squaring typical task-level correlations. In the discussion, we consider why these magnitudes differ and what they imply for interpreting the structure of individual differences in perceptual susceptibility.


Table~\ref{tab:CFA-Visual-Illusions} presents the results of the model, and Figure~\ref{fig:efhbm} visualizes key components of the variance decomposition. The first five factors capture variance specific to differences between versions of each illusion. The key result is that the sixth factor loads consistently across all ten tasks and explains over four times the variance of the seventh factor. This pattern strongly supports the conclusion that a single general factor underlies susceptibility to visual illusions. A projection of each task version onto the two general factors is shown in Figure~\ref{fig:efhbm}A. However, unique variance between the two versions of each illusion remains relevant (Figure~\ref{fig:efhbm}B), and unique variance attributable to specific versions of each illusion remains substantial (Figure~\ref{fig:efhbm}C). The general factor alone accounts for approximately 23.3\% of the total variance, which is substantially greater than the 5\% expected from squaring typical task-version-level correlations. In the discussion, we consider why these differences occur and how to interpret them in drawing scientific conclusions.


